% Sidechain Compressor parameter sweep results table.
% Requires packages already in main.tex's preamble: booktabs, siunitx.
% ALSO requires colortbl for \cellcolor (used for the grey parameter-column
% headers and the green/yellow heatmap columns below) - main.tex currently
% loads plain \usepackage{xcolor} with no [table] option, which does NOT
% pull in \cellcolor on its own. Either add \usepackage{colortbl} after
% xcolor, or change the xcolor line to \usepackage[table]{xcolor}.
%
% Data source: analysis/compute_folder_export_metrics.py, run against
% "UnMasking Exports/" (9 runs, Construction Scene, Dialogue priority class,
% Advanced/Summed-bus duck mode). Margin/STOI gain = processed - unprocessed.
% Ranked by mean margin gain, best to worst.
%
% Unprocessed baseline is the true sum of the 5 raw stems (Basic mode,
% Unmask off), not the scene's separately-bounced reference blend - the
% blend was found to differ from the raw stem sum by ~8% RMS overall (most
% likely a DAW export inconsistency), which risked contaminating the
% baseline with an artifact unrelated to ducking.
%
% Margin is measured on a band padded 0.85 octaves inward from Dialogue's
% nominal 400Hz-7kHz range (i.e. ~721-3883Hz), not the full nominal range -
% see compute_metrics.py's MARGIN_MEASUREMENT_PAD_OCTAVES. Reason: this
% script's own analysis bandpass filter, applied to the rendered output,
% interacts with Advanced mode's internal LR4 crossover (a different filter
% shape) near the shared nominal edges, inflating Basic mode's measured
% margin gain relative to Advanced's by a spurious ~0.9-1.0dB at the
% unpadded edges - confirmed via the engine's own reported gain that Basic
% and Advanced apply bit-identical gain within the class's core band when
% settings match, so that gap was a measurement artifact, not a real
% difference. 0.85 octaves of padding closes it to within 0.13dB on every
% one of these 9 runs (was ~0.9-1.0dB unpadded). This is a real reduction
% in measured bandwidth, accepted as a deliberate tradeoff.
%
% COLOR CODING: the five settings columns (#, Thr., Ratio, Knee, Atk.,
% Rel.) get grey header backgrounds - these are the manually-chosen sweep
% inputs, not measured outputs. Mean Margin Gain and STOI Gain are instead
% shaded per-CELL (not just the header), green and yellow respectively,
% intensity scaled linearly between this table's own min and max in each
% column (i.e. relative to these 9 rows only, not any absolute reference) -
% brighter/more saturated = higher gain. Because \cellcolor needs a plain
% number as its argument, not something siunitx's S-column type parses
% cleanly alongside, those two columns are plain right-aligned (r) rather
% than S columns here - each row's value is still written with consistent
% decimal precision by hand, so the visual alignment is the same either way.

\definecolor{paramgrey}{gray}{0.85}

\begin{table}[h]
\centering
\footnotesize
\caption{Sidechain Compressor parameter sweep: in-band masking-margin and STOI
gain of Advanced (Summed-bus) ducking over the unprocessed baseline
(Construction Scene, Dialogue priority class, measured on
\SI{721}{\hertz}--\SI{3883}{\hertz} - see note below). Ranked by mean
margin gain. Grey headers mark the manually-chosen sweep settings; the two
result columns are shaded per cell, more intense = higher gain.}
\label{tab:param-sweep}
\begin{tabular}{@{}c r r r r r r r@{}}
\toprule
\cellcolor{paramgrey}\textbf{\#} & \cellcolor{paramgrey}\textbf{Thr.} (\si{\decibel}) & \cellcolor{paramgrey}\textbf{Ratio} & \cellcolor{paramgrey}\textbf{Knee} (\si{\decibel}) & \cellcolor{paramgrey}\textbf{Atk.} (\si{\milli\second}) & \cellcolor{paramgrey}\textbf{Rel.} (\si{\milli\second}) & \textbf{Mean Margin Gain} (\si{\decibel}) & \textbf{STOI Gain} \\
\midrule
1 & -36 & 4.0 & 20 & 5  & 400 & \cellcolor{green!70}+8.51 & \cellcolor{yellow!68}+0.121 \\
2 & -48 & 1.5 & 12 & 5  & 400 & \cellcolor{green!68}+8.19 & \cellcolor{yellow!70}+0.124 \\
3 & -48 & 1.5 & 12 & 5  & 120 & \cellcolor{green!58}+6.78 & \cellcolor{yellow!64}+0.113 \\
4 & -36 & 1.5 & 12 & 1  & 120 & \cellcolor{green!36}+3.74 & \cellcolor{yellow!42}+0.068 \\
5 & -36 & 1.5 & 12 & 5  & 400 & \cellcolor{green!35}+3.58 & \cellcolor{yellow!39}+0.062 \\
6 & -48 & 1.5 & 12 & 50 & 120 & \cellcolor{green!34}+3.52 & \cellcolor{yellow!41}+0.066 \\
7 & -24 & 4.0 & 20 & 5  & 400 & \cellcolor{green!21}+1.65 & \cellcolor{yellow!23}+0.029 \\
8 & -36 & 1.5 & 12 & 15 & 120 & \cellcolor{green!21}+1.60 & \cellcolor{yellow!25}+0.033 \\
9 & -24 & 4.0 & 6  & 5  & 400 & \cellcolor{green!15}+0.83 & \cellcolor{yellow!15}+0.013 \\
\bottomrule
\end{tabular}
\par\vspace{4pt}
\raggedright\footnotesize Measurement band is padded 0.85 octaves inward
from Dialogue's nominal \SI{400}{\hertz}--\SI{7}{\kilo\hertz} range, not
the full nominal range: this script's own analysis filter, applied to the
rendered output, interacts with Advanced mode's internal crossover filter
near the shared nominal edges, inflating Basic mode's measured margin gain
relative to Advanced's by a spurious $\sim$\SI{0.9}{\decibel} at the
unpadded edges even though the two modes apply bit-identical gain within
the class's core band. Padding closes this to within
\SI{0.13}{\decibel} on every run tested here -- see
\texttt{analysis/compute\_metrics.py}'s \texttt{MARGIN\_MEASUREMENT\_PAD\_OCTAVES}
for full methodology.
\end{table}
